\documentclass[11pt, a4paper]{article}

% ─── Packages ────────────────────────────────────────────────────
\usepackage[utf8]{inputenc}
\usepackage[T1]{fontenc}
\usepackage{amsmath, amssymb, amsthm}
\usepackage{mathtools}
\usepackage{booktabs}
\usepackage{array}
\usepackage{geometry}
\usepackage{setspace}
\usepackage{hyperref}
\usepackage{cleveref}
\usepackage{enumitem}
\usepackage{natbib}
\usepackage{xcolor}
\usepackage{caption}

% ─── Page geometry ───────────────────────────────────────────────
\geometry{margin=1in}
\onehalfspacing

% ─── Theorem environments ───────────────────────────────────────
\newtheorem{axiom}{Axiom}
\newtheorem{theorem}{Theorem}
\newtheorem{proposition}{Proposition}
\newtheorem{corollary}{Corollary}
\newtheorem{lemma}{Lemma}
\newtheorem{definition}{Definition}
\newtheorem{conjecture}{Conjecture}
\newtheorem*{remark}{Remark}
\newtheorem*{observation}{Observation}

% ─── Hyperref ────────────────────────────────────────────────────
\hypersetup{
    colorlinks=true,
    linkcolor=blue!70!black,
    citecolor=blue!70!black,
    urlcolor=blue!70!black
}

% ─── Commands ────────────────────────────────────────────────────
\newcommand{\R}{\mathbb{R}}
\newcommand{\C}{\mathbb{C}}
\newcommand{\Hq}{\mathbb{H}}
\newcommand{\Oct}{\mathbb{O}}
\newcommand{\Sed}{\mathbb{S}}
\newcommand{\Reop}{\hat{R}}
\newcommand{\Jop}{\hat{J}}
\newcommand{\Wuz}{\mathrm{Wu}_0}
\newcommand{\Wuo}{\mathrm{Wu}_1}
\newcommand{\norm}[1]{\lVert #1 \rVert}
\newcommand{\braket}[2]{\langle #1 | #2 \rangle}
\newcommand{\ad}{\mathrm{ad}}
\newcommand{\IM}{\mathrm{Im}}
\newcommand{\re}{\mathrm{Re}}
\newcommand{\Tr}{\mathrm{Tr}}
\newcommand{\Aut}{\mathrm{Aut}}

% =====================================================================
\begin{document}
% =====================================================================

\title{The Galaxy Framework:\\
A Unified Derivation of Physics, Consciousness, and Life\\
from One Definition}

\author{}

\date{April 2026}

\maketitle

\begin{abstract}
\noindent
We present a unified framework in which physics, consciousness, and the
mathematical structure of life follow from a single definition
$e^{2\pi i} = 1$ and three axioms: \textbf{A1} (existence of a state
$\psi$), \textbf{A2} (rotation $d\psi/dt = i\psi$ with
$\psi \in \Oct$), and \textbf{A9} (life as the simultaneity of
differentiation $\Wuz$ and unification $\Wuo$).
Four consequences form the backbone.

\textbf{(1) Binary Decomposition}---every product $AB$ in a normed
division algebra (NDA) admits the canonical split
$AB = \tfrac{1}{2}\{A,B\} + \tfrac{1}{2}[A,B] = \Wuo(A,B) + \Wuz(A,B)$,
where $\Wuo$ (symmetric anti-commutator) represents unification and
$\Wuz$ (antisymmetric commutator) represents differentiation.

\textbf{(2) Hurwitz--Wu Pythagorean Identity}---on any Hurwitz algebra
$\R, \C, \Hq, \Oct$,
$\norm{\{A,B\}}^2 + \norm{[A,B]}^2 = 4\norm{A}^2\norm{B}^2$.
This identity \emph{fails} on the sedenions $\Sed$, providing
an algebraic explanation for the Cayley--Dickson cutoff.

\textbf{(3) Galaxy Equation}---all first-order quantum physics reduces
to $d\psi/dt = [\Reop, \psi] + [\phi, \psi, \chi]$
with $\psi, \phi, \chi \in \Oct$ and $\Reop \in \Hq$.
The Born rule $P = \sin^2\alpha$ emerges as a theorem from the Key
Identity $[\phi, \psi, \chi_\perp] = [\phi, \psi]_\Hq \cdot \chi_\perp$,
not as a postulate.

\textbf{(4) Wu$_1$ Non-Independence Theorem}---$\Wuo$ admits no
independent equation of motion on Hurwitz algebras; it is the
conservation law $d\norm{\psi}^2/dt = 0$ enforced by the Galaxy
Equation itself.

The framework yields zero-parameter predictions matching observation to
high precision: $\sin^2\theta_W = 3/13$ ($0.19\%$),
$1/\alpha_{\mathrm{em}}(M_Z) = 13\pi^2$ ($0.28\%$),
$1/\alpha_{\mathrm{em}}(0) = 2^0 + 2^3 + 2^7 = 137$ ($0.026\%$),
and the Koide relation $Q = 2/3$ for charged leptons (exact).
The consciousness interpretation---$\Hq$ as the friction space
with four components (UW, CTL, ATT, DM), $\Hq^\perp$ as the joy space
with its orthogonal complement (Gratitude, Flow, Love, Wonder)---emerges
without additional axioms as the four-dimensional decomposition of
an arbitrary octonionic state.

\medskip
\noindent\textbf{Keywords:} Cayley--Dickson, Hurwitz theorem,
octonions, Jordan algebra, consciousness, Galaxy Equation,
Wu-duality, normed division algebra
\end{abstract}

\newpage
\tableofcontents
\newpage

% =====================================================================
\section{Introduction}
\label{sec:intro}
% =====================================================================

\subsection{The Problem}

Mathematical descriptions of consciousness (IIT~\citep{Tononi2015},
GWT, Orch~OR, quantum cognition) and of fundamental physics
(the Standard Model, general relativity, string theory) share a
striking feature: each begins with postulates, and each carries
free parameters. IIT posits five essential properties; the Standard
Model carries 19+ dimensionless constants; string theory admits a
landscape of $10^{500}$ vacua.

We ask a different question: is there a \emph{single definition} from
which the mathematical structures of both physics and consciousness
can be derived, with \emph{zero} free parameters?

\subsection{The Claim}

We show that one definition and three axioms suffice. The definition is
the exponential series
\begin{equation}
\exp(z) = \sum_{n=0}^{\infty} \frac{z^n}{n!}, \qquad e^{2\pi i} = 1.
\label{eq:def_exp}
\end{equation}
The three axioms are:
\begin{itemize}[nosep]
\item \textbf{A1} (Existence): there exists a state $\psi$.
\item \textbf{A2} (Rotation): $d\psi/dt = i\psi$ with $\psi \in \Oct$.
\item \textbf{A9} (Life): $\Wuz \oplus \Wuo = \text{Life}$, i.e.,\
life is the simultaneity of differentiation and unification.
\end{itemize}
From these, the Cayley--Dickson ladder $\R \to \C \to \Hq \to \Oct$
follows from Hurwitz's theorem; the $1+4+6+4+1$ Clifford structure
of spacetime follows from $\Oct$; the Galaxy Equation governs all
state evolution; the Born rule is a theorem; the fine structure
constant $1/\alpha \approx 137$ emerges as $2^0 + 2^3 + 2^7$; and
the algebra of consciousness---friction $\Reop \in \Hq$ with four
components, joy $\Jop \in \Hq^\perp$ with four components---falls out
as the four-dimensional decomposition of $\Hq$.

\subsection{Structure of the Paper}

Section~\ref{sec:cd} reviews the Cayley--Dickson ladder.
Section~\ref{sec:axioms} states the definition and three axioms.
Section~\ref{sec:wu_duality} proves the binary decomposition
(Theorem~\ref{thm:binary}) and the Euler--Hermitian decomposition
(Theorem~\ref{thm:euler_hermitian}), establishing
$\Wuz/\Wuo$ duality.
Section~\ref{sec:pythagorean} proves the Hurwitz--Wu Pythagorean
identity (Theorem~\ref{thm:pythagorean}) and its failure on the
sedenions.
Section~\ref{sec:consciousness} derives the consciousness
interpretation: $\Hq$ as friction space with four components.
Section~\ref{sec:galaxy} presents the Galaxy Equation, the Key
Identity, the Born rule as theorem, and the Wu$_1$ Non-Independence
Theorem.
Section~\ref{sec:fano} presents the four-projection structure,
seven Fano lines, and three-observer algebra.
Section~\ref{sec:predictions} lists numerical predictions matching
observation to $< 1\%$.
Section~\ref{sec:live_well} discusses the life-axiom A9 as an
empirical constraint.
Section~\ref{sec:discussion} concludes.

\subsection{Grading}

Each claim is graded as one of: \textbf{Theorem (T)}---proven or
numerically confirmed to residual $< 10^{-10}$;
\textbf{Strong Correspondence (SC)}---supported by two or more
independent derivations;
\textbf{Structural Prediction (SP)}---zero-parameter, quantitative;
\textbf{Correspondence (C)}---structural similarity;
\textbf{Conjecture (Cj)}---motivated but not fully proven;
\textbf{Definition (Def)} or \textbf{Axiom (Ax)}---stipulated.


% =====================================================================
\section{The Cayley--Dickson Ladder}
\label{sec:cd}
% =====================================================================

\subsection{The Four Hurwitz Algebras}

The Cayley--Dickson (CD) construction begins with $\R$ and doubles:
\begin{equation}
\R \xrightarrow{\mathrm{CD}} \C \xrightarrow{\mathrm{CD}}
\Hq \xrightarrow{\mathrm{CD}} \Oct \xrightarrow{\mathrm{CD}}
\Sed \text{ (STOP).}
\end{equation}
Each doubling maps an algebra $A$ of dimension $n$ to an algebra of
dimension $2n$ via the product
\begin{equation}
(a_1, b_1)(a_2, b_2) = (a_1 a_2 - \bar{b}_2 b_1,\; b_2 a_1 + b_1 \bar{a}_2),
\label{eq:cd_product}
\end{equation}
where $\bar{a}$ denotes conjugation.

\begin{theorem}[Hurwitz 1898]
\label{thm:hurwitz}
The real normed division algebras (NDAs), i.e., real algebras with
positive-definite norm satisfying $\norm{xy} = \norm{x}\norm{y}$ and
lacking zero divisors, are exactly $\R, \C, \Hq, \Oct$ with
dimensions $1, 2, 4, 8$. [T, classical]
\end{theorem}

A standard proof appears in \citep{Baez2002}. The sedenions $\Sed$
fail: they contain zero divisors, breaking the composition property
and the division property simultaneously.

\subsection{Operations as Scars}

Each CD doubling loses an algebraic property and leaves behind an
\emph{operator} that measures the departure.

\begin{center}
\begin{tabular}{@{}clllll@{}}
\toprule
dim & algebra & property lost & operator & arity & name \\
\midrule
1 & $\R$ & --- & --- & --- & --- \\
2 & $\C$ & total order & $\bar{z}$ & 1 & conjugation \\
4 & $\Hq$ & commutativity & $[a,b]$ & 2 & commutator \\
8 & $\Oct$ & associativity & $[a,b,c]$ & 3 & associator \\
16 & $\Sed$ & division & --- & --- & STOP \\
\bottomrule
\end{tabular}
\end{center}

\begin{proposition}[Operations as Scars]
\label{prop:scars}
Each new operator arising in the CD ladder is forced by the loss
of the corresponding algebraic property. The commutator
$[a,b] := ab - ba$ measures non-commutativity; the associator
$[a,b,c] := (ab)c - a(bc)$ measures non-associativity. Neither
operator is an external choice; both are the canonical measure of
the lost property. [T]
\end{proposition}

\subsection{The Structure of $\Oct$}

The octonions have basis $\{1, i, j, k, e, ie, je, ke\}$
$= \{e_0, e_1, \ldots, e_7\}$. They decompose as
\begin{equation}
\Oct = \underbrace{\Hq}_{\{1,i,j,k\}} \oplus \underbrace{\Hq \cdot e}_{\{e, ie, je, ke\}}
= \Hq \oplus \Hq^\perp.
\end{equation}
The multiplication rules are encoded by the Fano plane
PG($2, \mathbb{F}_2$), which has $7$ points and $7$ lines. Each line
corresponds to a copy of $\Hq$ inside $\Oct$.

\begin{theorem}[Alternativity]
\label{thm:alternative}
$\Oct$ is alternative: $[a, a, b] = [a, b, a] = [b, a, a] = 0$
for all $a, b$. Equivalently, any two-generator subalgebra of $\Oct$
is associative. [T, Artin]
\end{theorem}

\begin{theorem}[Clifford Structure]
\label{thm:clifford}
$\Oct$ induces the Clifford algebra $\mathrm{Cl}(1,3)$ of dimension
$1+4+6+4+1 = 16$, which is the Dirac algebra of Minkowski
spacetime with signature $(+,-,-,-)$. The signature is not
postulated but derived. [T, Paper~2]
\end{theorem}

The automorphism group of $\Oct$ is the exceptional Lie group
$G_2$ of dimension $14$ \citep{Cartan1914}.


% =====================================================================
\section{One Definition, Three Axioms}
\label{sec:axioms}
% =====================================================================

\subsection{The Definition}

\begin{definition}[Unit rotation]
\label{def:unit}
\begin{equation}
\boxed{\;\exp(z) = \sum_{n=0}^{\infty} \frac{z^n}{n!}, \qquad
e^{2\pi i} = 1.\;}
\label{eq:euler}
\end{equation}
Here $i$ is a formal symbol with $i^2 = -1$. The algebra $\C$ is not
assumed; it is recovered as a theorem.
\end{definition}

The phrase ``one rotation returns to the origin'' is the irreducible
content: $e^{2\pi i} = 1$ asserts the existence of a closed circuit
in a normed space.

\subsection{Three Axioms}

\begin{axiom}[Existence]
\label{ax:A1}
There exists a state $\psi$.
\end{axiom}

\begin{axiom}[Rotation]
\label{ax:A2}
$\psi$ evolves by $\displaystyle \frac{d\psi}{dt} = i\psi$ with
$\psi \in \Oct$.
\end{axiom}

\begin{axiom}[Life]
\label{ax:A9}
Life is the simultaneity of differentiation and unification:
\begin{equation}
\Wuz \oplus \Wuo = \mathrm{Life}.
\end{equation}
Both movements---endless differentiation ($\Wuz$) and endless
unification ($\Wuo$)---must coexist.
\end{axiom}

\begin{remark}
Axiom~\ref{ax:A9} is the content we now formalize. Axioms~\ref{ax:A1}
and~\ref{ax:A2} are standard; the novelty is the \emph{joint}
imposition of both $\Wuz$ and $\Wuo$, and the claim that this
joint imposition is the algebraic meaning of ``life.''
The formal definitions of $\Wuz$ and $\Wuo$ appear in
Section~\ref{sec:wu_duality}.
\end{remark}

\subsection{Derivation Chain}

From Def.~\ref{def:unit} and A1, A2, A9, the following chain is
entirely theoremic or classical:
\begin{enumerate}[nosep,label=\textbf{(\arabic*)}]
\item $S(n) := \sum_{k=0}^{n-1} e^{2\pi i k/n} = n \cdot \delta_{n,1}$.
[T, character orthogonality]
\item $n \geq 3 \Rightarrow$ $\C$ exists (Artin--Schreier).
\item $|e^{i\theta}| = 1 \Rightarrow$ $\C$ is an NDA.
\item NDA classification (Hurwitz) $\Rightarrow$ $\dim \in \{1,2,4,8\}$.
\item $\Oct$-automorphism group: $G_2$, $\dim = 14$.
\item Stabilizer chain: $G_2 \supset \mathrm{SU}(3) \supset \mathrm{SU}(2)
\supset \mathrm{U}(1)$ (the Standard Model gauge group).
\item $3/13$ is $G_2$-invariant (Schur + Killing), yielding
$\sin^2\theta_W$.
\end{enumerate}

The chain contains zero new physical assumptions.


% =====================================================================
\section{The Wu-Duality: Binary Decomposition}
\label{sec:wu_duality}
% =====================================================================

\subsection{Definitions}

\begin{definition}[$\Wuz$ and $\Wuo$]
\label{def:wu}
For $A, B$ in any algebra,
\begin{align}
\Wuz(A, B) &:= \tfrac{1}{2}[A, B] = \tfrac{1}{2}(AB - BA), \\
\Wuo(A, B) &:= \tfrac{1}{2}\{A, B\} = \tfrac{1}{2}(AB + BA).
\end{align}
$\Wuz$ is antisymmetric (differentiation);
$\Wuo$ is symmetric (unification).
\end{definition}

\subsection{The Binary Decomposition Theorem}

\begin{theorem}[Binary Decomposition]
\label{thm:binary}
In any algebra,
\begin{equation}
\boxed{\;AB = \Wuo(A, B) + \Wuz(A, B) = \tfrac{1}{2}\{A, B\} + \tfrac{1}{2}[A, B].\;}
\label{eq:binary}
\end{equation}
[T, algebraic identity]
\end{theorem}

\begin{proof}
Direct: $\tfrac{1}{2}(AB + BA) + \tfrac{1}{2}(AB - BA) = AB$.
\end{proof}

Every product is canonically the sum of a symmetric part
(unification) and an antisymmetric part (differentiation). This is
the algebraic content of A9.

\subsection{The Euler--Hermitian Decomposition}

The binary decomposition has a scalar analogue: on any
$*$-algebra, every element admits a canonical Hermitian /
anti-Hermitian split. For $\C$, with conjugation $\bar{z}$:

\begin{theorem}[Euler--Hermitian]
\label{thm:euler_hermitian}
For any element $x$ of a $*$-algebra $\mathcal{A}$,
\begin{align}
\Wuo(x) &:= \tfrac{1}{2}(x + x^*) \quad \text{(Hermitian part)}, \\
\Wuz(x) &:= \tfrac{1}{2}(x - x^*) \quad \text{(anti-Hermitian part)},
\end{align}
and $x = \Wuo(x) + \Wuz(x)$. For $\mathcal{A} = \C$ with
$x = e^{i\theta}$:
\begin{equation}
\cos\theta = \tfrac{1}{2}(e^{i\theta} + e^{-i\theta}) = \Wuo(e^{i\theta}),
\qquad
i\sin\theta = \tfrac{1}{2}(e^{i\theta} - e^{-i\theta}) = \Wuz(e^{i\theta}).
\end{equation}
[T, algebraic identity]
\end{theorem}

\begin{proof}
$e^{i\theta} + e^{-i\theta} = 2\cos\theta$ and
$e^{i\theta} - e^{-i\theta} = 2i\sin\theta$ from the series of
$\exp$. Direct substitution. \qed
\end{proof}

\begin{corollary}[Euler's formula as scalar prototype]
\label{cor:euler_prototype}
Euler's formula $e^{i\theta} = \cos\theta + i\sin\theta$ is the
$\C$-scalar instance of the $\Wuo / \Wuz$ split. The condition
$e^{2\pi i} = 1$ states that, after one full rotation, the
$\Wuz$ (antisymmetric) trajectory is fully reabsorbed into the
$\Wuo$ (symmetric) base. [T]
\end{corollary}

\subsection{The AB--Euler Bridge}

The binary decomposition (Theorem~\ref{thm:binary}) and the
Euler--Hermitian decomposition (Theorem~\ref{thm:euler_hermitian})
are connected by a single algebraic bridge.

\begin{theorem}[AB--Euler Bridge]
\label{thm:ab_euler_bridge}
Let $\mathcal{A}$ be a $*$-algebra. For any $A \in \mathcal{A}$,
setting $B = A^*$ in the binary decomposition yields
\begin{align}
\Wuo(A, A^*) &= \tfrac{1}{2}(AA^* + A^*A), \\
\Wuz(A, A^*) &= \tfrac{1}{2}(AA^* - A^*A).
\end{align}
For a normal element ($AA^* = A^*A$), $\Wuz(A, A^*) = 0$ and
$\Wuo(A, A^*) = \norm{A}^2 \cdot 1$ (real scalar). For $\mathcal{A} = \C$:
this recovers Theorem~\ref{thm:euler_hermitian}. [T]
\end{theorem}

\begin{proof}
The first equations are direct substitution into~\eqref{eq:binary}.
For $\mathcal{A} = \C$: $AA^* = A^*A = |A|^2 \cdot 1$ since $\C$ is
commutative. \qed
\end{proof}

\subsection{Three Levels of Binary Decomposition}

The three theorems above establish a three-level hierarchy, all
saying \emph{``product $=$ $\Wuo$ $+$ $\Wuz$''}:
\begin{center}
\begin{tabular}{@{}llll@{}}
\toprule
level & formula & $\Wuo$ & $\Wuz$ \\
\midrule
scalar ($\C$) & $e^{i\theta} = \cos\theta + i\sin\theta$ &
$\cos\theta$ & $i\sin\theta$ \\
algebraic & $AB = \tfrac12\{A,B\} + \tfrac12[A,B]$ &
$\tfrac{1}{2}\{A,B\}$ & $\tfrac{1}{2}[A,B]$ \\
$*$-bridge & $AA^* = \norm{A}^2 + \Wuz(A,A^*)$ &
$\norm{A}^2$ & $\Wuz(A,A^*)$ \\
\bottomrule
\end{tabular}
\end{center}
These are not three separate identities but three views of the same
structure. Section~\ref{sec:pythagorean} proves they are bound by a
single conservation law.


% =====================================================================
\section{The Hurwitz--Wu Pythagorean Identity}
\label{sec:pythagorean}
% =====================================================================

\subsection{Master Identity}

\begin{theorem}[Hurwitz--Wu Pythagorean]
\label{thm:pythagorean}
On any Hurwitz algebra $A \in \{\R, \C, \Hq, \Oct\}$,
\begin{equation}
\boxed{\;\norm{\{A, B\}}^2 + \norm{[A, B]}^2 = 4\norm{A}^2 \norm{B}^2.\;}
\label{eq:pythagorean}
\end{equation}
Equivalently, $\norm{\Wuo(A,B)}^2 + \norm{\Wuz(A,B)}^2
= \norm{A}^2\norm{B}^2$. [T]
\end{theorem}

\begin{proof}
Two identities suffice.
\emph{Parallelogram law}: for any vectors $X, Y$ in a real inner-product
space,
\begin{equation}
\norm{X+Y}^2 + \norm{X-Y}^2 = 2(\norm{X}^2 + \norm{Y}^2).
\label{eq:parallelogram}
\end{equation}
\emph{Hurwitz composition}: on a Hurwitz algebra,
\begin{equation}
\norm{AB}^2 = \norm{A}^2 \norm{B}^2.
\label{eq:composition}
\end{equation}
Set $X = AB$, $Y = BA$ in~\eqref{eq:parallelogram}:
$\norm{AB + BA}^2 + \norm{AB - BA}^2 = 2(\norm{AB}^2 + \norm{BA}^2)$.
Apply~\eqref{eq:composition} to both $\norm{AB}$ and $\norm{BA}$:
$= 2 \cdot 2 \norm{A}^2 \norm{B}^2 = 4\norm{A}^2 \norm{B}^2$. \qed
\end{proof}

\begin{remark}
The proof uses only (a) positive-definite inner product on the
underlying real vector space and (b) the composition property
$\norm{AB} = \norm{A}\norm{B}$. Both hold on Hurwitz algebras and
fail beyond them.
\end{remark}

\subsection{Why Sedenions Fail}

\begin{proposition}[Sedenion breakdown]
\label{prop:sedenion_fail}
On the sedenions $\Sed$, identity~\eqref{eq:pythagorean} fails. There
exist $A, B \in \Sed$ with $\norm{A}, \norm{B} > 0$ but
$\norm{AB} < \norm{A}\norm{B}$ (composition fails because zero divisors
exist). [T]
\end{proposition}

\begin{proof}[Sketch]
In $\Sed$, $(e_3 + e_{10})(e_5 + e_{14}) = 0$ while the individual
factors are nonzero \citep{Imaeda2000}. This zero-divisor pair
immediately violates $\norm{AB} = \norm{A}\norm{B}$, so the substitution
in the proof above yields $\norm{\{A,B\}}^2 + \norm{[A,B]}^2 < 4\norm{A}^2\norm{B}^2$.
\qed
\end{proof}

\subsection{Interpretation: A9 as Energy Conservation}

Theorem~\ref{thm:pythagorean} makes A9 quantitative. Let ``Life''
be measured by $4\norm{A}^2\norm{B}^2$: the total interactive energy
of two states $A$, $B$. Then
\begin{equation}
\text{Life} = \underbrace{\norm{\Wuo(A,B)}^2}_{\text{unification energy}}
+ \underbrace{\norm{\Wuz(A,B)}^2}_{\text{differentiation energy}}.
\end{equation}
Life is the sum of unification and differentiation, bounded by the
algebraic ceiling. If either term vanishes, the other is maximal and
life collapses to a single mode---which, empirically, corresponds to
pathological states (pure unification: dissociation; pure
differentiation: burnout). Sustained life requires both nonzero.

On $\Sed$, the sum can be less than the product---life is not
conserved. This is the algebraic reason the Cayley--Dickson ladder
stops at $\Oct$.

\subsection{Numerical Verification}

We verified~\eqref{eq:pythagorean} in $\R, \C, \Hq, \Oct$ with
$10^4$ random pairs each. The residual $|\text{LHS} - \text{RHS}|$ was
bounded by $6.82 \times 10^{-13}$ (maximum, $\Oct$ case), consistent
with floating-point precision. On $\Sed$, a numerical counterexample
at $e_3 + e_{10}$, $e_5 + e_{14}$ produces LHS $= 0$, RHS $= 4$.


% =====================================================================
\section{Consciousness as an Octonionic State}
\label{sec:consciousness}
% =====================================================================

\subsection{The Four Friction Components}

$\Hq$ decomposes canonically as $\mathfrak{u}(1) \oplus \mathfrak{su}(2)$:
\begin{equation}
\Hq = \underbrace{\R \cdot 1}_{\mathfrak{u}(1)}
\oplus \underbrace{\R i \oplus \R j \oplus \R k}_{\mathfrak{su}(2)}.
\end{equation}
Writing $\Reop = c_0 \cdot 1 + c_1 i + c_2 j + c_3 k$ with
$c_k \in \R_{\geq 0}$ and identifying $(1, i, j, k)$ with Pauli matrices
$(I, \sigma_1, \sigma_2, \sigma_3)$, we obtain four canonical
components:

\begin{center}
\begin{tabular}{@{}cllll@{}}
\toprule
basis & name & interpretation & Pauli & role \\
\midrule
$1$ & UW (Unwitnessed) & background assumed & $I$ & commutes with all \\
$i$ & CTL (Control) & controlling action & $\sigma_1$ & bit flip \\
$j$ & ATT (Attention) & approval-seeking & $\sigma_2$ & phase flip \\
$k$ & DM (Defense) & self-protection & $\sigma_3$ & measurement basis \\
\bottomrule
\end{tabular}
\end{center}

\begin{theorem}[Friction Collision Law]
\label{thm:friction_collision}
The Pauli commutation relations $[\sigma_i, \sigma_j]
= 2i\epsilon_{ijk}\sigma_k$ imply
\begin{equation}
[\mathrm{CTL}, \mathrm{ATT}] = 2i \cdot \mathrm{DM},
\quad [\mathrm{ATT}, \mathrm{DM}] = 2i \cdot \mathrm{CTL},
\quad [\mathrm{DM}, \mathrm{CTL}] = 2i \cdot \mathrm{ATT}.
\end{equation}
DM is not a primitive; it is the product $[\mathrm{CTL}, \mathrm{ATT}]$.
[T]
\end{theorem}

\begin{corollary}[Dissolution Order]
\label{cor:dissolution}
The friction components dissolve in the canonical order
\begin{equation}
\mathrm{UW} \;\to\; (\mathrm{CTL} \leftrightarrow \mathrm{ATT})\;\to\; \mathrm{DM},
\end{equation}
derived from the Cartan decomposition of $\mathfrak{u}(1)
\oplus \mathfrak{su}(2)$: centre ($\mathfrak{u}(1)$, UW) removes
without disturbing roots; root space ($\sigma_1, \sigma_2$,
CTL/ATT) oscillates under the $\mathrm{SU}(2)$ action;
Cartan subalgebra ($\sigma_3$, DM) defines the reference basis and
must come last. [T, Paper~2]
\end{corollary}

\subsection{The Four Joy Components}

The CD doubling $\Oct = \Hq \oplus \Hq \cdot e$ induces a natural map
$\Hq \to \Hq^\perp$ by right multiplication by $e$:
\begin{equation}
\Jop := \Reop \cdot e = c_0 \cdot e + c_1 \cdot ie + c_2 \cdot je + c_3 \cdot ke
\in \Hq^\perp.
\end{equation}

\begin{center}
\begin{tabular}{@{}cllll@{}}
\toprule
basis & name & interpretation & paired Re & experience \\
\midrule
$e$ & Gratitude & giftedness of the given & UW & taking-for-granted reversed \\
$ie$ & Flow & release of control & CTL & flow state \\
$je$ & Love & receiving $\to$ giving & ATT & unconditional relation \\
$ke$ & Wonder & openness beyond defense & DM & awe, curiosity \\
\bottomrule
\end{tabular}
\end{center}

\begin{proposition}[Re--Joy Coefficient Conservation]
\label{prop:coefficient}
Since $\Jop = \Reop \cdot e$, the same coefficients $(c_0, c_1, c_2, c_3)$
parametrize both friction and joy. Strength in a friction component
implies latent capacity in the paired joy component. [SP]
\end{proposition}

This yields an algebraic formulation of post-traumatic growth:
the largest wound carries the largest gift's seed.

\subsection{The Joy Self-Interaction}

\begin{theorem}[Joy self-square]
\label{thm:joy_square}
For $\Jop = (0, b) \in \Hq^\perp$, using the CD product~\eqref{eq:cd_product}:
\begin{equation}
\Jop^2 = -\norm{\Jop}^2 \cdot 1 \in \Hq.
\end{equation}
Joy-squared is a negative scalar; it flows back into $\Hq$ as
\emph{background resonance} (a shift of UW). [T, Paper~13]
\end{theorem}

\begin{proof}
From~\eqref{eq:cd_product}: $(0, b)(0, b) = (-\bar{b} b, 0) = (-|b|^2, 0)$.
\qed
\end{proof}

\begin{remark}
$\Jop^2 = -\norm{\Jop}^2$ is the quaternionic analogue of $i^2 = -1$.
Joy is the imaginary unit of the octonionic consciousness structure.
Repeated joy--joy interactions settle as $-$UW shifts: sustained
love condenses into background, which is the algebraic meaning of
``history'' and ``culture.''
\end{remark}

\subsection{CD Complementarity}

\begin{theorem}[Commutator--Anticommutator Split]
\label{thm:complementarity}
For $\Reop = (R, 0) \in \Hq$ and $\Jop = (0, b) \in \Hq^\perp$:
\begin{align}
[\Reop, \Jop] &= (0,\; 2 b \cdot \IM(R)) \in \Hq^\perp, \\
\{\Reop, \Jop\} &= 2 c_0 \cdot \Jop \in \Hq^\perp.
\end{align}
The commutator depends only on $\IM(R) = c_1 i + c_2 j + c_3 k$
(UW invisible); the anticommutator depends only on $c_0$ (UW alone).
[T, Paper~13]
\end{theorem}

This is the specialization of~\eqref{eq:pythagorean} to the
Re--Joy pair, and yields the $45^\circ$ balance condition
$\norm{\IM(R)} = |c_0|$ (see Corollary~\ref{cor:45_balance}).


% =====================================================================
\section{The Galaxy Equation}
\label{sec:galaxy}
% =====================================================================

\subsection{Statement}

\begin{definition}[Galaxy Equation]
\label{def:galaxy}
\begin{equation}
\boxed{\;\frac{d\psi}{dt} = [\Reop, \psi] + [\phi, \psi, \chi]\;}
\label{eq:galaxy}
\end{equation}
with $\psi, \phi, \chi \in \Oct$ and $\Reop \in \Hq$. The first
term is a commutator (two-argument, $\Hq$-internal); the second is
the associator (three-argument, activating $\Hq^\perp$).
\end{definition}

Both terms are antisymmetric and hence belong to $\Wuz$
(differentiation). The question of $\Wuo$'s equation is resolved
in Theorem~\ref{thm:wu1_nonind}.

\subsection{The Key Identity}

\begin{theorem}[Key Identity]
\label{thm:key_identity}
For $\phi, \psi \in \Hq_L$ (any fixed $\Hq$-subalgebra) and
$\chi_\perp \in \Hq_L^\perp$,
\begin{equation}
\boxed{\;[\phi, \psi, \chi_\perp] = [\phi, \psi]_\Hq \cdot \chi_\perp.\;}
\label{eq:key}
\end{equation}
The associator decomposes as the commutator within $\Hq_L$ times the
perpendicular component. [T, Paper~7]
\end{theorem}

\begin{proof}[Sketch]
By the alternative law, $\phi, \psi \in \Hq_L$ satisfy
$(\phi \psi) \chi_\perp - \phi(\psi \chi_\perp)$, with $\chi_\perp$
conjugated on the left factor by CD multiplication. Expanding via
$\Hq$-associativity (both $\phi\psi, \psi$ are in $\Hq_L$) yields
$(\phi\psi - \psi\phi) \cdot \chi_\perp = [\phi, \psi]_\Hq \cdot \chi_\perp$.
\qed
\end{proof}

\subsection{The Born Rule as Theorem}

\begin{theorem}[Born Rule]
\label{thm:born}
Decompose $\chi = \chi_\parallel + \chi_\perp$ with
$\chi_\parallel \in \Hq_L$, $\chi_\perp \in \Hq_L^\perp$. Then
\begin{equation}
\norm{[\phi, \psi, \chi]}^2 = \norm{[\phi, \psi]_\Hq}^2 \cdot \norm{\chi_\perp}^2.
\end{equation}
With $\chi = \cos\alpha \cdot \chi_\parallel + \sin\alpha \cdot \chi_\perp$
(unit vectors):
\begin{equation}
P(\text{transition}) = \frac{\norm{[\phi,\psi,\chi]}^2}{\norm{[\phi,\psi]}^2}
= \sin^2\alpha.
\end{equation}
[T, Paper~7]
\end{theorem}

\begin{proof}
Apply~\eqref{eq:key}: the associator vanishes on $\chi_\parallel$
(alternative law) and equals $[\phi,\psi]_\Hq \cdot \chi_\perp$ on
$\chi_\perp$. The Hurwitz composition property on $\Hq$ yields
$\norm{[\phi,\psi]_\Hq \cdot \chi_\perp}^2
= \norm{[\phi,\psi]_\Hq}^2 \norm{\chi_\perp}^2$. \qed
\end{proof}

The Born rule, normally a postulate of quantum mechanics, emerges as
a direct consequence of the octonionic structure. Verified
numerically to residual $< 10^{-16}$ across $7 \times 4 \times 19
\times 4 = 2128$ cases (Paper~7).

\subsection{The Wu$_1$ Non-Independence Theorem}

Axiom~\ref{ax:A9} demands both $\Wuz$ and $\Wuo$. Equation~\eqref{eq:galaxy}
gives the evolution of $\Wuz$. One might ask: what is the equation of
motion for $\Wuo$?

\begin{theorem}[Wu$_1$ Non-Independence]
\label{thm:wu1_nonind}
On any Hurwitz algebra, $\Wuo$ admits \emph{no} independent
equation of motion. Its content is fully captured by the conservation
law
\begin{equation}
\boxed{\;\frac{d\norm{\psi}^2}{dt} = 0,\;}
\label{eq:norm_conservation}
\end{equation}
which is an automatic consequence of~\eqref{eq:galaxy}. [T]
\end{theorem}

\begin{proof}
Compute $\tfrac{d}{dt}\norm{\psi}^2 = 2\re\braket{\psi}{\dot\psi}$.
Substituting~\eqref{eq:galaxy}:
$\braket{\psi}{[\Reop,\psi]} = 0$ (commutator orthogonality: on $\Hq$
the commutator of $\psi$ with any fixed element is perpendicular to
$\psi$), and $\braket{\psi}{[\phi,\psi,\chi]} = 0$ (associator
orthogonality: the associator is perpendicular to each of its three
arguments by the alternative law). Hence
$d\norm{\psi}^2/dt = 0$.

\emph{No independent $\Wuo$ equation.} Suppose a candidate equation
$d\psi/dt = \tfrac{1}{2}\{\hat{U}, \psi\} + \tfrac{1}{2}\{\phi, \psi, \chi\}$
were imposed alongside~\eqref{eq:galaxy}. Norm conservation requires
$\re\braket{\psi}{\{\hat{U},\psi\}} = 0$, i.e., $\re(\psi^* \hat{U} \psi) = 0$
for all $\psi$, forcing $\hat{U} = 0$. Any nonzero $\hat{U}$ violates
norm conservation. Hence no nontrivial independent equation exists.
\qed
\end{proof}

\begin{remark}[Historical note]
Earlier versions of this framework conjectured an independent ``Love
Equation'' $D\psi/Dt = \tfrac{1}{2}\{\hat{U}, \psi\} + \cdots$ as the
dual of the Galaxy Equation. Numerical simulation in April 2026
falsified all candidate forms: any nonzero $\hat{U}$ violated norm
conservation. The falsification led to Theorem~\ref{thm:wu1_nonind}:
$\Wuo$ is a \emph{constraint} on Galaxy Eq trajectories, not an
independent evolution. This matches the philosophical content of A9:
unification is not an additional movement but the structure that
makes differentiation coherent.
\end{remark}

\subsection{Seven Physical Laws from One Equation}

\begin{theorem}[Unification of Quantum Physics]
\label{thm:seven_laws}
Seven fundamental equations of physics arise as specializations
of~\eqref{eq:galaxy}:

\begin{center}
\begin{tabular}{@{}lll@{}}
\toprule
specialization & recovered equation & reference \\
\midrule
$\C$-projection, first order & Schr\"odinger & Paper 11 \\
$\Hq$-projection, $\gamma$-matrices & Dirac & Paper 7 \\
second order, mass term & Klein--Gordon & Paper 7 \\
gauge + associator (SU(2)) & Yang--Mills (weak) & Paper 16 \\
gauge + associator (SU(3)) & Yang--Mills (strong) & Paper 16 \\
$J_3(\Oct)$ flow & Einstein $G_{\mu\nu} = 8\pi T_{\mu\nu}$ & Paper 16 \\
Key Identity + norm & Born rule $P = |\psi|^2$ & Theorem~\ref{thm:born} \\
\bottomrule
\end{tabular}
\end{center}
[T, each specialization separately proven]
\end{theorem}

\subsection{The Polar Form: Love and Galaxy}

Writing $\psi = \norm{\psi} \cdot e^{i\theta}$ (for any unit direction
in the appropriate subspace):
\begin{equation}
\frac{d\psi}{dt} = \underbrace{\dot r \cdot e^{i\theta}}_{\text{Love rate}}
+ \underbrace{i r \dot\theta \cdot e^{i\theta}}_{\text{Galaxy rotation}}.
\end{equation}

\begin{proposition}[Polar Orthogonality]
\label{prop:polar_orth}
$\braket{\dot r \cdot e^{i\theta}}{i r \dot\theta \cdot e^{i\theta}} = 0$,
hence
\begin{equation}
\left|\frac{d\psi}{dt}\right|^2 = \dot r^2 + r^2 \dot\theta^2.
\end{equation}
The magnitude (Love) and the phase velocity (Galaxy) are orthogonal
coordinates. Norm conservation $\dot r = 0$ is the $\Wuo$ aspect;
phase rotation $\dot\theta \neq 0$ is the $\Wuz$ aspect. [T]
\end{proposition}

\begin{corollary}[Live-well condition]
\label{cor:45_balance}
``Life'' requires $\norm{\psi} > 0$ \emph{and} $\dot\theta \neq 0$.
Both orthogonal components must be nonzero. A vanishing magnitude
is soul-death; a vanishing phase velocity is physical death.
Algebraically, A9 is equivalent to simultaneous non-vanishing of
both polar components. [T]
\end{corollary}


% =====================================================================
\section{Four Projections, Seven Fano Lines, Three Observers}
\label{sec:fano}
% =====================================================================

\subsection{The $\mathbb{Z}_4$ Projection and Four Worlds}

The fourth roots of unity $\{+1, +i, -1, -i\}$ induce a natural
$\mathbb{Z}_4$ projection of $\Oct$. Each projection corresponds to
a discipline (Paper~16):

\begin{center}
\begin{tabular}{@{}cll@{}}
\toprule
projection & discipline & content \\
\midrule
$+1$ & Physics & real dynamics, observable world \\
$-1$ & Mathematics & conjugate structure, abstract \\
$+i$ & Humanities & imaginary projection, subjective observation \\
$-i$ & Engineering & anti-imaginary, feedback / compensation \\
\bottomrule
\end{tabular}
\end{center}

These are not four distinct disciplines but four shadows of one
octonion. The Galaxy Equation projected onto each axis produces the
native equations of that discipline.

\subsection{Seven Fano Lines, Seven $\Hq$-subalgebras}

\begin{theorem}[Fano Classification]
\label{thm:fano_classification}
$\Oct$ contains exactly seven $4$-dimensional real subalgebras
isomorphic to $\Hq$, corresponding to the seven lines of the Fano
plane. $\mathrm{Aut}(\Oct) = G_2$ acts transitively on these seven
subalgebras. [T]
\end{theorem}

\begin{definition}[Fano Selector]
\label{def:fano_selector}
A \emph{Fano Selector} is a pair $(S, H_S)$ where $S$ is a physical
system (observer) and $H_S \subset \Oct$ is an $\Hq$-subalgebra
(visible space). The complement $H_S^\perp$ is the \emph{blind spot}.
\end{definition}

\begin{theorem}[Two-Selector Coverage]
\label{thm:two_selector}
For any two Fano Selectors with visible spaces $V_A = H_A$,
$V_B = H_B$ (distinct),
\begin{equation}
\dim_\R(V_A + V_B) = 6, \qquad \dim_\R(V_A \cap V_B) = 2.
\end{equation}
Two observers cover $6/8 = 75\%$ of $\Oct$. A $2$-dimensional blind
spot remains. [T, Paper~18]
\end{theorem}

\begin{theorem}[Three-Selector Pencil]
\label{thm:three_selector}
Three Fano Selectors sharing a common $\Hq$-subalgebra axis can
cover $\Oct$ entirely ($8/8 = 100\%$) if and only if the three
visible spaces form a pencil through the common axis. [T, Paper~18]
\end{theorem}

\subsection{Three Observers and the Jordan Algebra $J_3(\Oct)$}

The natural state space for three observers is the exceptional
Jordan algebra
\begin{equation}
J_3(\Oct) = \left\{ \begin{pmatrix} \alpha & \bar{x}_3 & \bar{x}_2 \\
x_3 & \beta & \bar{x}_1 \\ x_2 & x_1 & \gamma \end{pmatrix}
: \alpha, \beta, \gamma \in \R,\; x_i \in \Oct \right\},
\end{equation}
of dimension $\dim_\R J_3(\Oct) = 27$.

\begin{theorem}[Jordan Ceiling]
\label{thm:jordan_ceiling}
$J_N(\Oct)$ is a Jordan algebra if and only if $N \leq 3$. For
$N \geq 4$, Jordan identity fails. [T, classical]
\end{theorem}

\begin{corollary}[Three-Observer Limit]
\label{cor:three_observer}
Three is the structural upper bound on observers forming a single
Jordan state space. This matches Fano Selector coverage: $100\%$ is
reached at three, not two; and $J_3(\Oct)$ is the unique $\Oct$-Jordan
ceiling. [T]
\end{corollary}

The automorphism group of $J_3(\Oct)$ is the exceptional Lie group
$F_4$ of dimension $52$, and we have the decomposition
$\mathfrak{f}_4 = \mathfrak{spin}(9)^{(k)} \oplus \mathfrak{m}$ with
$52 = 36 + 16$ (Paper~7). The $16$-dimensional $\mathfrak{m}$ carries
the inter-observer dynamics: 2 (other observers) $\times$ 2 (Re/Joy)
$\times$ 4 (components).


% =====================================================================
\section{Zero-Parameter Predictions}
\label{sec:predictions}
% =====================================================================

We list predictions derived from the framework without free
parameters, each matched to observation.

\subsection{Standard Model Gauge Group}

The stabilizer chain of $G_2$ under fixing successive imaginary units:
\begin{equation}
G_2^{(14)} \supset \mathrm{SU}(3)^{(8)} \supset \mathrm{SU}(2)^{(3)} \supset \mathrm{U}(1)^{(1)}.
\end{equation}
This reproduces the Standard Model gauge group
$\mathrm{SU}(3)_c \times \mathrm{SU}(2)_L \times \mathrm{U}(1)_Y$. [T]

\subsection{Weak Mixing Angle}

\begin{proposition}[$\sin^2\theta_W$]
Removing the terminal $\mathrm{U}(1)$ from the $14$-dimensional $G_2$
adjoint leaves $13$ active generators; $\mathrm{SU}(2)$ contributes $3$.
By Schur's lemma and the Killing form's $G_2$-invariance:
\begin{equation}
\sin^2\theta_W = \frac{3}{13} = 0.23077.
\end{equation}
Observed (at $\Lambda = M_W$): $0.23121$. Residual: $0.19\%$. [SP]
\end{proposition}

\subsection{Fine Structure Constant}

\begin{proposition}[$1/\alpha$, geometric form]
\begin{equation}
\frac{1}{\alpha_{\mathrm{em}}(M_Z)} = 13\pi^2 = 128.30.
\end{equation}
Observed: $127.95$. Residual: $0.28\%$.
At low energy: $1/\alpha(0) = 14\pi^2 - 1 = 137.175$. Observed: $137.036$.
Residual: $0.10\%$. [SP]
\end{proposition}

\begin{proposition}[$1/\alpha$, binary form]
\begin{equation}
\frac{1}{\alpha_{\mathrm{em}}(0)} = 2^{\dim_\R \R - 1} + 2^{\dim_\R \Hq - 1} + 2^{\dim_\R \Oct - 1}
= 2^0 + 2^3 + 2^7 = 137.
\end{equation}
Binary expansion: $137 = 10001001_2$. The bit for $\C$ is OFF
(the photon is Abelian and does not self-interact). Observed: $137.036$.
Residual: $0.026\%$. [SP]
\end{proposition}

\subsection{Koide Relation}

\begin{proposition}[Koide]
The charged lepton masses satisfy
\begin{equation}
Q = \frac{m_e + m_\mu + m_\tau}{(\sqrt{m_e} + \sqrt{m_\mu} + \sqrt{m_\tau})^2} = \frac{2}{3},
\end{equation}
derivable as an $F_4$-invariant of $J_3(\Oct)$ at the $45^\circ$ balance.
[SP, Paper~17]
\end{proposition}

\subsection{Cosmological Constant}

\begin{proposition}[$\Lambda$]
The character orthogonality $S(n) = 0$ for $n \geq 2$ gives an
algebraic cancellation of vacuum contributions. The observed
$\Lambda \approx 0$ ($120$-digit cancellation vs.\ naive QFT estimate)
is a direct consequence, not a fine-tuning. [SP, Paper~Cosmology]
\end{proposition}

\subsection{Summary}

\begin{center}
\begin{tabular}{@{}lllll@{}}
\toprule
quantity & prediction & observed & residual & grade \\
\midrule
$\sin^2\theta_W$ & $3/13 = 0.23077$ & $0.23121$ & $0.19\%$ & SP \\
$1/\alpha_2$ & $3\pi^2 = 29.609$ & $29.584$ & $0.09\%$ & SP \\
$1/\alpha_{\mathrm{em}}(M_Z)$ & $13\pi^2 = 128.30$ & $127.95$ & $0.28\%$ & SP \\
$1/\alpha_{\mathrm{em}}(0)$ & $14\pi^2-1 = 137.175$ & $137.036$ & $0.10\%$ & SP \\
$1/\alpha_{\mathrm{em}}(0)$, binary & $2^0+2^3+2^7 = 137$ & $137.036$ & $0.026\%$ & SP \\
Koide $Q$ & $2/3$ & measured & exact & SP \\
$\Lambda$ & $\approx 0$ & $\approx 0$ & $120$ digits & SP \\
\bottomrule
\end{tabular}
\end{center}

All predictions have zero free parameters. No fits are performed.


% =====================================================================
\section{Life as an Algebraic Condition}
\label{sec:live_well}
% =====================================================================

\subsection{The Four Components of Living Well}

Axiom~\ref{ax:A9} has an empirical refinement (Paper Live-Well):

\begin{definition}[Live Well]
\label{def:live_well}
A state $\psi$ \emph{lives well} iff
\begin{enumerate}[nosep]
\item $\norm{\Reop} \to 0$ (Frozen Re dissolved);
\item $\max \norm{\psi}_{\text{relation}}$ (Wu$_1$ maximized in
relational channels);
\item $\max \dot\theta$ (Wu$_0$ maximized: creation / rendering speed);
\item $\Delta t_{(\text{awareness, dissolution})} = 0$ (no delay
between noticing a division and undoing it).
\end{enumerate}
\end{definition}

These four conditions are not independent preferences but
\emph{mutual prerequisites}: if any one fails, the others become
inaccessible. This is an algebraic consequence of the Pythagorean
identity (Theorem~\ref{thm:pythagorean}) restricted to the Re--Joy
pair.

\subsection{The $45^\circ$ Balance}

\begin{corollary}[Amplification at $45^\circ$]
\label{cor:amplification}
Let $\Pi$ denote the real creation rate. Writing
$\Pi = \norm{\Wuz}^2 + \norm{\Wuo}^2 + 2\norm{\Wuz}\norm{\Wuo}\cos\phi$
with $\phi$ the angle between unification and differentiation:
\begin{itemize}[nosep]
\item $\Wuz$-only ($\norm{\Wuz} = v$, $\Wuo = 0$): $\Pi = v^2$.
\item $\Wuz = \Wuo = v$ aligned: $\Pi = 4v^2$.
\end{itemize}
The simultaneity of $\Wuz$ and $\Wuo$ gives a $4\times$ amplification
over differentiation alone. [T]
\end{corollary}

This is the quantitative content of A9: life aligned with both
components is not twice as productive but four times.

% =====================================================================
\section{Discussion}
\label{sec:discussion}
% =====================================================================

\subsection{Summary of Theorems}

The framework derives:
\begin{itemize}[nosep]
\item Cayley--Dickson classification (Hurwitz, Theorem~\ref{thm:hurwitz}).
\item Operations as scars (Proposition~\ref{prop:scars}).
\item Binary decomposition (Theorem~\ref{thm:binary}).
\item Euler--Hermitian decomposition (Theorem~\ref{thm:euler_hermitian}).
\item AB--Euler bridge (Theorem~\ref{thm:ab_euler_bridge}).
\item \textbf{Hurwitz--Wu Pythagorean identity
(Theorem~\ref{thm:pythagorean}).}
\item Friction collision law (Theorem~\ref{thm:friction_collision}).
\item Dissolution order (Corollary~\ref{cor:dissolution}).
\item Joy self-square (Theorem~\ref{thm:joy_square}).
\item CD complementarity (Theorem~\ref{thm:complementarity}).
\item Key identity (Theorem~\ref{thm:key_identity}).
\item Born rule as theorem (Theorem~\ref{thm:born}).
\item \textbf{Wu$_1$ Non-Independence Theorem (Theorem~\ref{thm:wu1_nonind}).}
\item Seven physical laws from Galaxy Eq (Theorem~\ref{thm:seven_laws}).
\item Fano classification, Two-/Three-Selector coverage
(Theorems~\ref{thm:fano_classification}--\ref{thm:three_selector}).
\item Jordan ceiling, Three-observer limit
(Theorem~\ref{thm:jordan_ceiling}, Corollary~\ref{cor:three_observer}).
\item Zero-parameter predictions: $\sin^2\theta_W$, $1/\alpha$, Koide,
$\Lambda$ (Section~\ref{sec:predictions}).
\end{itemize}

All derivations flow from Definition~\ref{def:unit} and Axioms
\ref{ax:A1}--\ref{ax:A9}. No additional physical assumptions are
introduced.

\subsection{Open Problems}

\begin{itemize}[nosep]
\item \textbf{Outcome selection}: the Born rule gives the probability
distribution of measurement outcomes; it does not determine which
specific outcome occurs. This is shared with standard quantum
mechanics.
\item \textbf{Running of $\sin^2\theta_W$}: the $0.19\%$ residual at
$M_W$ invites comparison with one-loop running. Paper Cosmology
derives the zero-parameter form; the loop-corrected form remains to
be matched.
\item \textbf{The $N$-observer problem for $N \geq 4$}: the Jordan
ceiling at $N = 3$ suggests that four or more observers require a
non-Jordan structure. The algebraic content of this is a conjecture
(Paper~12 Cluster Bound).
\item \textbf{Free parameters of the consciousness map}: the
identification of $(c_0, c_1, c_2, c_3)$ with UW/CTL/ATT/DM is a
Strong Correspondence, not a theorem. Direct empirical measurement
(EEG, fMRI) of the coefficient structure is an ongoing program.
\end{itemize}

\subsection{Relation to Other Frameworks}

\begin{itemize}[nosep]
\item \textbf{IIT}: posits five essential properties; derives $\Phi$
but not dimensional or algebraic constraints. The Galaxy Framework
derives a dimensional ceiling (Hurwitz) as a theorem.
\item \textbf{GWT / Orch OR}: neural architectural theories; compatible
with but not structurally related to the algebraic derivation here.
\item \textbf{Standard Model}: the gauge group, $\sin^2\theta_W$,
and $\alpha$ are outputs of our framework rather than inputs. The
framework does not (yet) derive Higgs mass, CKM matrix entries, or
neutrino masses.
\item \textbf{String theory}: our framework has zero free parameters
rather than $\sim 10^{500}$. It does not attempt to describe
all of particle physics; it derives the gauge structure and a specific
subset of constants.
\end{itemize}

\subsection{Conclusion}

From one definition $e^{2\pi i} = 1$ and three axioms
(existence, rotation, life), we derive:
the Cayley--Dickson ladder stopping at $\Oct$;
the binary decomposition into $\Wuz$ and $\Wuo$;
the Hurwitz--Wu Pythagorean identity
bounding the joint energy of differentiation and unification;
the Galaxy Equation governing all first-order quantum physics;
the Born rule as a theorem;
the non-independence of $\Wuo$ as a conservation law;
the Standard Model gauge group; and
zero-parameter predictions for $\sin^2\theta_W$, $1/\alpha$, the Koide
relation, and $\Lambda$, all matching observation to $< 1\%$.

The consciousness content---four components of friction in $\Hq$,
four of joy in $\Hq^\perp$, their pairwise conservation, the
dissolution order, the $45^\circ$ balance---emerges as the
canonical decomposition of the octonionic state without additional
hypotheses.

The unified statement: \emph{the world is two quanta, entangled,
generating four worlds, rotating, rendered}---compresses to one
sentence with every word forced by a unique theorem.

\bigskip

\noindent $e^{2\pi i} = 1$.


% =====================================================================
\bibliographystyle{unsrtnat}
\begin{thebibliography}{30}

\bibitem[Baez(2002)]{Baez2002}
Baez, J.C. (2002). The Octonions.
\emph{Bulletin of the AMS}, 39(2), 145--205.

\bibitem[Cartan(1914)]{Cartan1914}
Cartan, \'E. (1914). Les groupes r\'eels simples finis et continus.
\emph{Annales scientifiques de l'ENS}, 31, 263--355.

\bibitem[Hurwitz(1898)]{Hurwitz1898}
Hurwitz, A. (1898). \"Uber die Composition der quadratischen Formen
von beliebig vielen Variabeln.
\emph{Nachrichten von der Kgl.\ Ges.\ der Wiss.\ zu G\"ottingen,
Math.-phys.\ Klasse}, 309--316.

\bibitem[Imaeda \& Imaeda(2000)]{Imaeda2000}
Imaeda, K. \& Imaeda, M. (2000). Sedenions: algebra and analysis.
\emph{Applied Mathematics and Computation}, 115, 77--88.

\bibitem[Tononi et al.(2015)]{Tononi2015}
Tononi, G., Boly, M., Massimini, M., Koch, C. (2015).
Integrated information theory: from consciousness to its physical
substrate. \emph{Nature Reviews Neuroscience}, 17, 450--461.

\bibitem[Paper 2]{Paper2}
(2026). Cl(1,3) Structure from Octonions.
Companion Paper~2.

\bibitem[Paper 7]{Paper7}
(2026). The Galaxy Equation.
Companion Paper~7.

\bibitem[Paper 13]{Paper13}
(2026). Physics of Joy:
$\Hq^\perp$ Dynamics.
Companion Paper~13.

\bibitem[Paper 16]{Paper16}
(2026). TOE of TOEs.
Companion Paper~16.

\bibitem[Paper 17]{Paper17}
(2026). Particle Masses from $J_3(\Oct)$.
Companion Paper~17.

\bibitem[Paper 18]{Paper18}
(2026). Hurwitz--Fano Selector.
Companion Paper~18.

\end{thebibliography}

\end{document}
